\ChapterPage{Página 6 \textbullet{} Transcripción}{El minuto perdido}
\noindent\textbf{VALE:} ¿Qué ocurrió cuando sonó la campana?\\
\textbf{BELL:} El pasadizo se llenó de luz. No era una luz brillante. Hacía
que cada sombra pareciera una puerta.\\
\textbf{VALE:} ¿Cuántos golpes?\\
\textbf{BELL:} Tres.\\
\textbf{VOSS:} Cuatro.\\
\textbf{VALE:} No están de acuerdo.\\
\textbf{VOSS:} No estamos de acuerdo en el orden. El cuarto llegó primero.

\begin{quote}
\textbf{VALE:} El reloj de la ciudad se detuvo a las 3:33.\\
\textbf{BELL:} Se detuvo porque oyó la campana.\\
\textbf{VALE:} Los relojes no oyen.\\
\textbf{BELL:} Este sí.
\end{quote}

\textbf{VALE:} ¿Qué fue de la carta?\\
\textbf{BELL:} El papel estaba en blanco cuando llegamos a la calle.\\
\textbf{VALE:} ¿Y la botella?\\
\textbf{BELL:} Seguía seca.\\
\textbf{VALE:} El registro muestra aquí una pausa de cincuenta y nueve
segundos.\\
\textbf{VOSS:} Su reloj perdió un minuto.\\
\textbf{VALE:} Mi reloj funciona bien.\\
\textbf{VOSS:} La marea también.

\medskip
\noindent\textit{En este punto la grabación contiene un tono grave. El
escribiente no anotó más preguntas. La transcripción termina con las palabras
«puerta abierta», escritas con una letra que no es la suya.}
